% Generated draft: requires mathematical and source review.
\documentclass[11pt]{article}
\usepackage[margin=25mm]{geometry}
\usepackage{fontspec}
\setmainfont{Latin Modern Roman}
\usepackage{amsmath,amssymb,mathtools,unicode-math}
\setmathfont{Latin Modern Math}
\usepackage{enumitem,needspace,xurl,hyperref}
\setlength{\parindent}{0pt}
\setlength{\parskip}{5pt}
\setlength{\emergencystretch}{3em}
\newcommand{\N}{\mathbb{N}}
\newcommand{\Z}{\mathbb{Z}}
\newcommand{\Q}{\mathbb{Q}}
\newcommand{\R}{\mathbb{R}}
\newcommand{\C}{\mathbb{C}}
\newcommand{\F}{\mathbb{F}}
\newcommand{\A}{\mathbb{A}}
\newcommand{\PP}{\mathbb{P}}
\newcommand{\E}{\mathbb{E}}
\newcommand{\eps}{\varepsilon}
\newcommand{\dd}{\,\mathrm{d}}
\begin{document}
\section*{UnsolvedMath ID: 30002037}
\textbf{Source problem code:} catalogue entry
\subsection{Definitions and mathematical statement}
Let \(X\) be a toric variety of dimension \(n\), let \(D\) be a divisor on \(X\), and let
\[
U\subseteq H^0(X,\mathcal O_X(D)),\qquad \dim U=n+2,
\]
be a basepoint-free subspace. The ground field, the class of admissible toric varieties, and the type of divisor are not specified in the supplied record. If \(\mathcal O_X(D)\) is an invertible sheaf, basepoint-freeness means that the evaluation map
\[
U\otimes\mathcal O_X\longrightarrow\mathcal O_X(D)
\]
is surjective. For a general divisor for which this sheaf is not invertible, the intended meaning requires clarification.

Under the invertible-sheaf interpretation, a basis \(s_0,\ldots,s_{n+1}\) of \(U\) defines the morphism
\[
\phi_U:X\longrightarrow\mathbb P^{n+1},\qquad
x\longmapsto[s_0(x):\cdots:s_{n+1}(x)].
\]
Intrinsically, its target parametrizes one-dimensional quotients of \(U\), equivalently lines in \(U^*\). The source writes this target as \(\mathbb P(U)\); that notation must be read with the corresponding projectivization convention.

Assume, as in the supplied corrected statement, that the image is a hypersurface. Denote the intended hypersurface by \(Y\). For implicitization one ordinarily considers the Zariski closure of the image, but the record does not distinguish the image from its closure or specify a scheme structure. Once these choices and a ground field \(k\) are fixed, an implicit equation means a nonzero homogeneous polynomial
\[
F\in k[T_0,\ldots,T_{n+1}]
\]
defining the chosen hypersurface, up to multiplication by an element of \(k^\times\).

\textbf{Requested construction and geometric description.} For the data above, develop an efficient method to compute an implicit equation of \(Y\) from an approximation complex, and describe how the multigraded syzygies of the ideal denoted by \(I_U\) control the geometry of \(Y\). If a specified graded coordinate ring \(S\) represents the sections \(s_i\) by homogeneous elements \(f_i\), a possible interpretation is
\[
I_U=(f_0,\ldots,f_{n+1})\subseteq S,
\]
with multigraded syzygies given by homogeneous relations \(\sum_i a_i f_i=0\) and their higher relations. The record does not specify \(S\), its grading, the representation of the sections, or the approximation complex; this interpretation is therefore conditional.

The original statement additionally requires the image to be a \emph{toric} hypersurface, whereas the supplied corrected statement requires only a hypersurface. These targets are retained as an unresolved difference. The meaning of efficiency and the geometric properties to be recovered from syzygies are also unspecified. No generic-finiteness or birationality hypothesis is supplied, and an equation-extraction procedure must state any additional hypotheses it needs. This is a research-program statement requiring clarification, not a fully specified theorem.
\subsection{Short English statement}
Given a toric variety and an \((n+2)\)-dimensional basepoint-free space of divisor sections whose image is a hypersurface, seek an efficient approximation-complex method for finding its equation and a description of how multigraded syzygies determine its geometry. The field, divisor assumptions, graded ring, approximation complex, meaning of efficiency, and intended geometric conclusions remain unspecified. The original and corrected records also differ on whether the image must itself be toric.
\subsection{Sources}
\begin{itemize}
\item Local registry: supplied UnsolvedMath source record, permanent ID \(30002037\), source problem code \texttt{catalogue entry}, title \emph{Toric Hypersurface Implicitization from Syzygies}; collection: Oberwolfach Reports Open Problems. No individual proposer is supplied; the recorded proposed year is 2012.
\item Supplied report citation: \emph{Toric Geometry (2013)}, Oberwolfach Reports, pp.\ 1231--1306, DOI \texttt{10.4171/owr/2012/21}, \url{https://doi.org/10.4171/owr/2012/21}. The citation's 2013 title date and its 2012 DOI component and proposed-year metadata are preserved without reconciliation.
\item Supplied literature source, labelled Botbol--Dickenstein: \url{https://doi.org/10.1016/j.jsc.2015.09.001}.
\item Supplied literature source, labelled Recent syzygy implicitization: \url{https://arxiv.org/abs/2606.01415}.
\item Supplied literature source, labelled Toric Geometry report: \url{https://ems.press/content/serial-article-files/46392}.
\item Literature assessment checked at \texttt{2026-08-21}: the supplied curation labels the problem partially solved and states that approximation-complex and syzygy methods cover broad surface cases and selected higher-dimensional cases, that recent work extends matrix constructions, and that arbitrary toric hypersurface implicitization remains an active general problem. This is supplied dated curation derived from OWR material, not an independently verified literature assessment.
\item The registry labels the statement \texttt{corrected\_verified} and reports that the official report supplies the toric-variety, divisor, subspace, dimension, and hypersurface hypotheses omitted or distorted in extraction. This verification claim is supplied metadata; the report and cited literature have not been independently inspected for this draft.
\end{itemize}
\end{document}
